\subsection{\textsf{Induction}}
\subsubsection{Weak induction}
Inductions are usually used when something must be proved for a large number of
objects, usually for every number of objects.
\begin{example}{}{}
	Prove that $n^2+n$ is even for every positive integer $n$.
\end{example}
\textcolor{color2}{\textbf{Solution. }} It is clearly true, because $n^2+n=n(n+1)$. Moreover, of course $n$ or $n+1$ must be even.
However, we can prove it using induction either as we can see in the step-by-step below.
\begin{enumerate}
	\item Base case. Suppose we wanna prove $P(n)$ is true for every $n$ real number. Therefore, the first step in induction
	      is proving that $P(n_0)$ is true. This is called the base case.

	      If we wanna prove $n^2+n$ is even, we need to show that $1^2+1=2$ is even first, which clearly is.
	\item Inductive step. Now, we suppose that $P(k)$ is true when $n_0<k<k+1$. This is called the inductive hypotheses. And
	      then, we prove that $P(k+1)$ is true using the hypotheses that $P(k)$ is true. If we manage to do this,
	      every $P(n)$ is going to be true.
\end{enumerate}
\subsubsection{Multiple base cases}
Now, say you want to prove the following using induction.
\begin{example}{}{}
	Suppose you have and unlimited supply of $5$ and $8$ cent stamps. Show you can use these to make any amount of postage greater than $27$ cents.
\end{example}
\textcolor{color2}{\textbf{Solution. }} The best way to solve this one is proving that $n=28,29,30,31,32$ can all be obtained using only out $5$ and $8$ cent stamps.
\begin{itemize}
	\item $28=5\cdot4 + 8\cdot1$
	\item $29=5\cdot1 + 8\cdot3$
	\item $30=5\cdot6 + 8\cdot0$
	\item $31=5\cdot3 + 8\cdot2$
	\item $32=5\cdot0+ 8\cdot4$
\end{itemize}
As we can see, instead of $1$ base case, we have $5$. Now, we must simply prove that $P(k+5)$ is true after supposing that
$P(k)$ is true. Nevertheless, it clearly is, because we must simply add a $5$ cent stamp to make it true.

You might have already noticed the generalization for this one. If you want to prove $P(n)$ is true, you must simply prove
$P(k+m)$ is true after supposing $P(k)$ is true and making all the base cases, that are $P(n_0),\dots, P(m)$.
\subsubsection{Strong Induction}
The strong induction, differently from the weak induction, uses multiple inductive hypotheses. To be more specific,
we assume $P(j)$ is true for every $j<n$ and try to prove that $P(n)$ is true based on our hypotheses.
\begin{example}{}{}\label{exa:strong}
	You and I are playing a game with raisins. We start with two non-empty piles of raisins
	on the table, and take turns making moves. A move consists of eating all the raisins in one pile and
	dividing the second into two non-empty (and not necessarily equal) piles. The game continues until
	one of us can’t move; that player loses. Show that if one of the piles starts out with an even number
	of raisins then the first player can win.
\end{example}
\textcolor{color2}{\textbf{Solution. }}
Let $a$ and $b$ denote the current number of raisins in each pile on the table. Say
\begin{itemize}
	\item $(a,b)=(1,1)$. Then, it is impossible to divide any pile into two non-empty piles. Therefore, the current player to move loses.
	\item $(a,b)=(2,1)$ or $(a,b)=(1,2)$. Then, it suffices to eat the $1$ raisin pile and divide the $2$ raisins pile. Therefore, the
	      current player to move wins.
\end{itemize}
We proceed by strong induction on $a+b$. Let's suppose the following is true for every $a+b<c$:
\begin{itemize}
	\item $a$ and $b$ are both odd. Hence, the current player to move loses.
	\item At least one of the piles has an even number of raisins. Hence, the current player to move wins.
\end{itemize}
Alright, so let's prove now that $a+b=c$ follows our hypotheses:
\begin{itemize}
	\item $a$ and $b$ are both odd. Hence, one of the piles will be divided into two piles of opposite parity, while the other one is eaten.
	      Thus, according to our hypotheses, the next player will win, since there is at least one pile filled with an even number of raisins.
	      Consequently, the current player loses.

	\item At least one of the piles has an even number of raisins. Hence, the current player can divide the pile with an even number of raisins
	      into a $1$ raisin pile and a $a-1$ or $b-1$ raisins pile (depending on which pile the current player ate). Therefore,
	      both of the resulting piles are odd. According to our hypotheses, the current player wins, because the total number of raisins is either
	      $a$ or $b$, and $a,b<a+b$.
\end{itemize}
Therefore, the first player has the winning strategy when at least one of the piles has an even number of raisins.
\subsubsection{Weak induction vs. Strong induction}
The biggest difference between those two, is that:
\begin{itemize}
	\item Weak induction.
	      \begin{itemize}
		      \item Base case: must prove that $P(0), \dots, P(m)$ is true. Usually, $m=0$. In this case,
		            we must simply prove $P(0)$ is true.
		      \item Inductive step: must assume $P(n)$ is true for every $0 \le n \le n+1$ in order to demonstrate that $P(n+1+m)$ is true.
	      \end{itemize}
	\item Strong induction.
	      \begin{itemize}
		      \item Base case: Base case: must prove that $P(0), \dots, P(m)$ is true. Usually, $m=0$. In this case,
		            we must simply prove $P(0)$ is true.
		      \item Inductive step: must assume $P(0), \dots, P(n-1)$ is true for every $0 \le n$ in order to demonstrate that $P(n)$ is true.
	      \end{itemize}
\end{itemize}

\subsubsection{Inverted induction}
This name was actually created by me. The inverted induction ain't really inverted, it is just simply different.

Suppose you want to prove that $P(n)$ is true for all $n\in \NN$.
\begin{itemize}
	\item Define a set $N$ containing all the $n$ that result in $P(n)$ false.
	\item By the Well-Ordering Principle, there must be a least elemenet $k\in S$ if $S$ is non-empty. This one is often called the minimal criminal or the smallest counterexample.
	\item Demonstrate by contradiction that $S$ must be empty.
\end{itemize}
Usually it is said that $n\in\NN$. However, we simply need $S$ to have a least number. Therefore, $S$ can't be $(0,1)$, for example, but can be
$[0,1)$, even though it is not natural and not finite.

To contradict the minimal criminal, we usually
\begin{itemize}
	\item show that every $n<k$ satisfies $P(n)$ and then show $P(k)$ is true either (basically like a strong induction).
	\item show that the smallest counterexample isn't actually the smallest, because there is an even smaller counterexample.
\end{itemize}

\begin{example}{}{}
	Prove $\sqrt{2}$ is irrational.
\end{example}
\textcolor{color2}{\textbf{Solution. }} For the sake of contradiction, say $\sqrt{2}$ is rational. Then there exists $p,q>0$ such that $\sqrt{2}=\frac{p}{q}$,
with $q$ the smallest possible postive denominator.
Hence,
\[
	\sqrt{2} = \frac{p}{q}\implies 2=\frac{p^2}{q^2} \iff 2q^2=p^2.
\]
Consequently, $p=2p'$, meaning that $p$ is even. Therefore
\[
	2q^2=4p'^2 \iff q^2=2p'^2.
\]
Consequently, $q=2q'$, meaning that $q$ is even.

\[
	\sqrt{2}=\frac{p}{q} \iff \sqrt{2}=\frac{2p'}{2q'} = \frac{p'}{q'}.
\]
Contradiction! We said that $q$ was the least denominator possible, but $q'<q$. Hence $q$ isn't the least.
\begin{example}{}{}
	Prove that every $n \geq 2$ has a prime factorization.
\end{example}
\textcolor{color2}{\textbf{Solution. }} By the sake of contradiction, say there exists a set $N$ containing every $n\ge 2$ that doesn't have a prime factorization.
Hence, by the WOP, there is a least element $n'$ contained in $N$.

If $n'$ were prime, it would have a prime factorization. Thus it isn't prime and $n'=ab$. Since $n'$ is the smallest number inside $N$, $a$ and $b$ must
have a prime factorization. Contradiction! $n'$ has a prime factorization that is the product of the prime factorizations of $a$ and $b$.

\subsubsection{Examples}
\begin{example}{}{}
	$P$ is true for all $n,m\ge1$. Prove
	\[
		\frac{m!}{0!} + \cdots + \frac{(m+n)!}{n!} = \frac{(m+n+1)!}{n!(m+1)}.
	\]
\end{example}
\textcolor{color2}{\textbf{Solution. }}
\begin{itemize}
	\item Base case: $n=1$ results in
	      \[
		      \frac{m!}{0!} + \frac{(m+1)!}{1!} = \frac{(m+2)!}{m+1},
	      \]
	      which is clearly true.
	\item Inductive step: Suppose $P(n)$ is true. Therefore
	      \[
		      P(n+1)= \frac{m!}{0!}+\cdots+\frac{(m+n+1)!}{(n+1)!}= \frac{(m+n+1)!}{(n+1)!} + \frac{(m+n+1)!}{n!(m+1)}
		      = \frac{(m+n+2)!}{(n+1)!(m+1)}
	      \]
	      can easily be proved by expanding $\frac{(m+n+1)!}{(n+1)!} + \frac{(m+n+1)!}{n!(m+1)}$.
\end{itemize}

\begin{insight}
	Inducting over $m$ is harder since using the fact that $P(m)$ is true isn't as easy as using the fact that $P(n)$ is true.
	This kinda of situation often happens, when there are two or more possibilities to induct over, but one of them is easier.
\end{insight}

\begin{example}{}{}
	Prove the AM-GM inequality.
\end{example}
\textcolor{color2}{\textbf{Solution. }}

\clearpage

\subsection{\textsf{Problems}}

\subsubsection{USAJMO 2019 P1}

\begin{problem}
There are $a+b$ bowls arranged in a row, numbered $1$ through $a+b$, where $a$
and $b$ are given positive integers. Initially, each of the first $a$
bowls contains an apple, and each of the last $b$ bowls contains a
pear. A legal move consists of moving an apple from bowl $i$ to bowl $i+1$
and a pear from bowl $j$ to bowl $j-1$, provided that the difference $i-j$ is
even. We permit multiple fruits in the same bowl at the same time. The goal is
to end up with the first $b$ bowls each containing a pear and the last $a$
bowls each containing an apple. Show that this is possible if and only if the
product $ab$ is even.
\end{problem}

We claim that, in order to reach the desired configuration, $a$ and $b$ cannot
both be odd; equivalently, $ab$ must be even. We will first show that an even
product $ab$ always allows us to reach the goal. Then we will prove that the
desired configuration is impossible when $ab$ is odd.

Let $i'$ be the apple originally in bowl $i$ and $j'$ be the pear
in $j$.

\begin{claim}
	If $i-j$ is even, then $i'$ can always be swapped with $j'$.
\end{claim}

\begin{proof}
	WLOG assume $i<j$. After $n$ legal moves, $i'$ and $j'$ are in bowls $i+n$ and
	$j-n$, respectively. Since $i-j$ is even, $n=\frac{j-i}{2}$ is an integer.
	Therefore, $i'$ and $j'$ meet at
	\[
		i+n=j-n=\frac{i+j}{2}.
	\]
	After that, $i'$ can move to $\frac{i+j}{2}+n$ to reach $j$ and $j'$ to
	$\frac{i+j}{2}-n$ to reach $i$.
\end{proof}

Now suppose $ab$ is even.

\begin{itemize}
	\item If $\min(a,b)=1$, say $a=1$, then $b$ must be even for $ab$ to be even.
	      Therefore the bowls $1$ and $a+b$ can be swapped, because $(a+b)-1$ is
	      even.
	\item If $\min(a,b)\ge 2$
	      \begin{itemize}
		      \item and $a+b$ is odd, $(a+b)-1$ is even, so the first and the last
		            bowls can be swapped. Therefore, $ab$ can be reduced to $(a-1)
			            (b-1)$ since the bowls $1$ and $a+b$ already have the desired
		            fruits, which is still even. Thus, by complete induction, $a+b$
		            odd works, since $(a-1)+(b-1)$ is smaller than $a+b$.
		      \item and $a+b$ is even, $(a+b)-2$ and $(a+b-1)-1$ are even. $a$ and $b$
		            are both even numbers, because if they were both odd numbers, $ab$
		            wouldn't be even. So the
		            second and the last bowls can be swapped, as the first and
		            penultimate bowls. Therefore, we can reduce $ab$ to $(a-1)(b-1)$,
		            which is still even, since the bowls $1$, $2$, $a+b-1$ and $a+b$
		            already have the desired fruits. Thus, by complete induction,
		            $a+b$ even works, since $(a-2)+(b-2)$ is smaller than $a+b$.
	      \end{itemize}
\end{itemize}
Hence, when $ab$ is even, we can reach the desired goal.

Now, by the sake of contradiction, say $ab$ is odd. Let $A$ be the amount of
apples in the first odd-numbered bowls and $B$ of pears in the last odd-numbered
$b$ bowls. We already know that $i$ and $j$ must have parity in order to $i-j$ be
even. Therefore $A$ and $B$ are invariants, because after a legal move, $i'$ and
$j'$ must keep having parity, so $A$ and $B$ are always going to each decrease by
$1$ or increase by $1$.

Since $A-B$ is invariant, the remainder must always be equal, regardless of the fruits'
order (apples first, then pears, or pears first, then apples). Since $a$ is odd,
but $a+b$ is even, $A=\frac{a+1}{2}$ and $B=\frac{b-1}{2}$ originally, but after
the swaps, $A=\frac{a-1}{2}$ and $B=\frac{b+1}{2}$. Consequently,
\[
	\frac{a+1}{2}-\frac{b-1}{2}=\frac{a-1}{2}-\frac{b+1}{2}\iff\frac{a-b+2}{2}=
	\frac{a-b-2}{2}\iff 2=-2
\]
Contradiction! Hence, we can reach the desired goal if and only if $ab$ is even.

\begin{insight}
	I am proud of myself because this is the first serious combinatorics problem I
	solved, even though there were a lot of mistakes in my original solution.

	\begin{itemize}
		\item Lack of rigorosity. I understand now that writing more to make the
		      solution more comprehensible is good. However, time matters a lot in
		      an olympiad context. Therefore, solving all the problems before actually
		      writing the solutions seems to be the best thing to do. This way, I can
		      keep track of my time better, understand if I can or not be that rigorous.
		      Another thing about writing solutions is being more "mathematician".
		      What I mean is that writing solutions mathematically is often faster
		      and more understandable.
		\item Learned about induction, an awesome tool to use when an iteration or
		      anything related must be demonstrated. In this problem, we needed to show
		      that an specific configuration could only be obtained when $ab$ were even.
		      However, notice how many iterations and steps there are here!
		\item Learned about invariants either. Can be used in a similar way. Usually
		      when something changes a lot, there is another thing that actually never
		      changes, which can be useful to find contradictions.
	\end{itemize}
\end{insight}

\clearpage

\subsubsection{USAJMO 2016 P4}
\begin{problem}
Find, with proof, the least integer $N$ such that if any $2016$ elements are
removed from the set ${1, 2,...,N}$, one can still find $2016$ distinct numbers
among the remaining elements with sum $N$.
\end{problem}

The least integer $N$ that satisfies the statement is
\[
	\sum^{4032}_{i=2017} i = 6049 \cdot 1008 = 6097392.
\]

Notice that if we form pairs of numbers from the set with equal sum, there will be
at least $3024$ such pairs. Even if at most $2016$ pairs are destroyed, there will
still remain at least $1008$ pairs with equal sum, i.e., $2016$ numbers. Each pair
consists of the $x$th number from the left and the $x$th number from the right.
This way, each pair has sum $6048+1$ as in $(1,6048), (2,6047), \dots, (3024,3025)$.
$6049 \cdot 1008 = 6097392$.

$N$ can't be less than $6,097,392$ because the least possible sum is
\[
	\sum^{2016}_{i=1} i=2033136.
\]
However, the first $2016$ numbers of the set can be removed, changing the least
possible sum to
\[
	\sum^{4032}_{i=2017} i = 6097392.
\]

\begin{insight}
	The most important idea that didn't came to my mind was forming pairs with equal
	sum, since I had already thought about lower bound $N \ge \sum^{2016}_{i=1}i$
	(not entirely correct, but almost there). It feels like the biggest problem
	was focusing more on doing something instead of thinking about what to do.
	Therefore, from now on it is really important to think better about what I am
	going to do before just doing.
\end{insight}

\clearpage

\subsubsection{Shortlist 2012 C1}

\begin{problem}
Several positive integers are written in a row. Iteratively, Alice chooses two adjacent numbers $x$ and $y$ such that $x>y$ and $x$ is to the left of $y$, and replaces the pair $(x,y)$ by either $(y+1,x)$ or $(x-1,x)$. Prove that she can perform only finitely many such iterations.
\end{problem}
\textcolor{color2}{\textbf{Solution. }}

\clearpage

